\chapter{Considerações Finais}
\label{chp:con}

\section{Síntese do Trabalho}
\label{sec:sintese}

Esta dissertação investigou a hipótese central de que a incorporação de conhecimento físico-químico diretamente na engenharia de atributos --- em vez de restringi-lo às funções de perda ou às arquiteturas de rede --- produz modelos de regressão supervisionada mais acurados, mais robustos e mais interpretáveis para a predição de parâmetros de fim de corrida em Fornos Básicos a Oxigênio (BOF). A motivação partiu de um diagnóstico conhecido na literatura: modelos puramente orientados a dados tendem a apresentar degradação significativa de desempenho quando aplicados a regiões operacionais fora da distribuição de treino, especialmente em processos siderúrgicos com dados escassos e ruidosos \cite{bae2020bof,willard2020integrating}.

A resposta metodológica proposta organiza-se em quatro etapas sequenciais. Primeiro, um conjunto de dados sintéticos de BOF foi gerado via simulação de Monte Carlo, com N=10.000 corridas fisicamente válidas, impondo restrições de balanço entálpico e estequiométrico às amostras \cite{buggineni2024synthetic,melo2024synthetic}. Segundo, os dados brutos foram expandidos com oito atributos físico-derivados --- demanda teórica de oxigênio, eficiência de sopro, entalpia acumulada, fração de massa oxidada, taxa de descarbonetação, basicidade da escória, capacidade de fusão da sucata e perda de ferro --- todos calculados a partir das equações de balanço descritas no Capítulo~\ref{chp:capitulo4}. Terceiro, três configurações experimentais foram definidas: Config~A (baseline, atributos brutos, conjunto completo de modelos otimizados via Optuna), Config~B (aprendizado profundo, atributos brutos, MLP com busca de arquitetura via Optuna) e Config~C (física informada, atributos brutos mais os oito físico-derivados, conjunto completo de modelos). Quarto, os modelos foram avaliados por MAE, RMSE e $R^2$ para três variáveis-alvo: temperatura final ($T_{\text{final}}$), teor de carbono final ($[\text{C}]_{\text{final}}$) e teor de fósforo final ($[\text{P}]_{\text{final}}$), adotando as tolerâncias industriais de $\text{MAE} \leq 10\,^\circ\text{C}$, $\text{MAE} \leq 0{,}020\,\%$ e $\text{MAE} \leq 0{,}003\,\%$, respectivamente, como critério de aprovação \cite{zhou2024bof,bae2020bof}.

A hipótese central foi parcialmente confirmada: os atributos físico-derivados produziram ganhos estatisticamente significativos para $[\text{C}]_\text{final}$ (10/11 modelos melhoram, $p < 0{,}05$) e $[\text{P}]_\text{final}$ (8/11 modelos melhoram), com seis modelos de $[\text{P}]_\text{final}$ passando a atender ao limiar industrial de $0{,}003\,\%$ contra apenas um na linha de base, mas não resolveram o gargalo estrutural de $T_\text{final}$ (nenhum modelo melhora significativamente), cujo erro é dominado por variância irredutível associada a grandezas estocásticas não observáveis ($\eta_\text{eff}$). O teste de extrapolação mostra que esses ganhos em $[\text{C}]_\text{final}$ e $[\text{P}]_\text{final}$ se mantêm sob perturbação moderada fora da faixa de calibração. A análise SHAP é consistente, em grande parte, com a relevância física esperada dos atributos: \texttt{o2\_efficiency} é a feature de maior importância para $[\text{C}]_\text{final}$ nos três modelos testados individualmente, e \texttt{slag\_basicity} para $[\text{P}]_\text{final}$ (embora com margem modesta sobre o teor bruto de fósforo do gusa) --- sem que isso prove, por si só, captura causal de estrutura física\,\cite{manojlovic2022eaf,xia2024dmpinn}.


\section{Principais Contribuições}
\label{sec:contribuicoes}

O trabalho entrega quatro contribuições concretas à interseção entre aprendizado de máquina e modelagem de processos siderúrgicos:

\textbf{C1 — Dataset sintético de BOF com validação física.} Um conjunto de 10.000 corridas sintéticas de BOF, amostrado dentro de faixas operacionais documentadas na literatura \cite{chattopadhyay2022hybrid,bae2020bof} e validado por consistência termodinâmica (erro de balanço entálpico $< 10\%$). O dataset supre a carência crítica de dados industriais abertos para pesquisa em predição de fim de corrida BOF, problema amplamente reconhecido na área \cite{ghalati2023review,zhou2024bof}. O código de geração e as constantes físicas utilizadas são disponibilizados integralmente no repositório do projeto\footnote{\url{https://github.com/devGvieira/bof-piml-dissertation}}, permitindo reprodução e extensão a outras faixas operacionais.

\textbf{C2 — Conjunto padronizado de oito atributos físico-derivados para BOF.} A sistematização das oito features calculadas a partir de balanços de massa e energia --- com equações, constantes termodinâmicas e referências bibliográficas explícitas para cada uma --- oferece um ponto de partida reproduzível para estudos futuros de engenharia de atributos em processos de refino. Ao contrário de abordagens que integram a física à função de perda \cite{raissi2019pinns,xia2024dmpinn}, a estratégia adotada é agnóstica em relação ao algoritmo de aprendizado e pode ser aplicada a qualquer modelo de regressão supervisionada.

\textbf{C3 — Comparação sistemática de três configurações experimentais em onze modelos e três variáveis-alvo.} O protocolo A×B×C isola dois fatores independentes: o ganho proveniente de atributos físicos (A×C) e o ganho proveniente de profundidade arquitetural (A×B), além de avaliar sua interação (B×C). Com 11 modelos (Ridge, Lasso, ElasticNet, SVR, KNN, RF, GBM, XGBoost, LightGBM, CatBoost, MLP) avaliados em três alvos, o experimento gera evidência quantitativa sobre qual fator --- física ou capacidade de representação --- contribui mais para a qualidade preditiva no domínio BOF. Protocolos similares de comparação sistemática em contextos siderúrgicos são escassos na literatura \cite{ghalati2023review,shao2024hybrid}.

\textbf{C4 — Análise SHAP com validação de consistência física.} A análise SHAP (SHapley Additive exPlanations) é aplicada não apenas como ferramenta de ranqueamento de importância de atributos, mas como mecanismo de validação qualitativa, sob um protocolo de falsificação pré-registrado (Capítulo~\ref{chp:capitulo1}): quatro das cinco hipóteses de coerência física testáveis (H1–H5) não foram falsificadas, o que é consistente com --- ainda que não prove --- a captura de estrutura física real pelos modelos, e não apenas de correlações espúrias \cite{manojlovic2022eaf,xia2024dmpinn}. A exceção (H5) e a fragilidade de H1 e H4 (Capítulo~\ref{chp:resultados}) revelam um limite metodológico específico: a média de importância SHAP entre modelos explicados por métodos distintos (TreeExplainer exato vs.\ KernelExplainer aproximado) pode ser dominada por um único modelo de escala bruta maior, mascarando divergência de sinal ou de posição entre os demais — um resultado negativo tratado como achado, não omitido. Esta perspectiva estende a aplicação de SHAP em siderurgia, onde seu uso é predominantemente descritivo na literatura revisada nesta dissertação \cite{ghalati2023review}.


\section{Limitações}
\label{sec:limitacoes}

A validade externa dos resultados desta pesquisa está condicionada a três limitações estruturais que devem ser consideradas na interpretação das conclusões. Ressalta-se que nenhum dado industrial proprietário foi empregado nesta pesquisa; todos os parâmetros físicos, faixas operacionais e constantes termodinâmicas derivam exclusivamente de literatura científica publicada\,\cite{cai2021heat,bae2020bof,chattopadhyay2022hybrid,xia2024dmpinn}.

Primeiramente, toda a experimentação é conduzida sobre dados sintéticos, gerados por um modelo físico simplificado que não inclui dinâmicas de escória, perfis temporais de sopro, variações de composição do sucateado ou efeitos de desgaste refratário. Comparações entre configurações AI e AI-física sobre dados sintéticos quantificam o ganho relativo da engenharia de atributos dentro do universo do modelo generativo, mas não permitem inferência direta sobre desempenho em plantas reais. A validação com dados industriais reais é, portanto, condição necessária para afirmar aplicabilidade operacional --- e permanece como trabalho futuro essencial \cite{buggineni2024synthetic,morales2025synthetic}.

Em segundo lugar, o modelo físico subjacente é estático e de estado final: prediz os valores de saída ao término da corrida a partir de variáveis de entrada iniciais, sem modelar a trajetória temporal do processo. Fenômenos como a dinâmica de formação de escória, o perfil de sopro ao longo do tempo e a evolução do banho metálico são ignorados. Modelos que incorporam perfis temporais --- como os abordados por \citeonline{shao2024hybrid} com dados de off-gas e por \citeonline{xia2024dmpinn} com PINNs dinâmicas --- podem capturar informações adicionais não disponíveis nesta formulação.

Finalmente, a generalização dos modelos treinados para plantas siderúrgicas distintas daquela que deu origem aos parâmetros físicos adotados pode exigir re-calibração das constantes, dos limites de amostragem e dos pesos do modelo físico. As tolerâncias industriais adotadas ($\text{MAE} \leq 10\,^\circ\text{C}$ para temperatura, $\leq 0{,}020\,\%$ para carbono e $\leq 0{,}003\,\%$ para fósforo) refletem condições de operação típicas documentadas na literatura \cite{zhou2024bof,bae2020bof}, mas podem diferir entre plantas com diferentes capacidades de convertedor, qualidades de sucata e fontes de ferro-gusa.


\section{Trabalhos Futuros}
\label{sec:trabalhos-futuros}

As limitações identificadas na seção anterior delineiam naturalmente as extensões mais relevantes desta pesquisa.

A prioridade imediata é a validação com dados reais de planta. A parceria com uma siderurgia nacional para acesso a históricos de corridas --- mesmo que anonimizados --- permitiria avaliar se os ganhos preditivos observados em dados sintéticos se transferem para o domínio real, e em que magnitude. A metodologia de engenharia de atributos físico-derivados é, por construção, transferível: as equações de balanço são universais, e apenas os parâmetros numéricos requerem calibração para cada instalação \cite{chattopadhyay2022hybrid,shao2024hybrid}.

Uma extensão natural do trabalho é a incorporação de perfis temporais de sopro e dados de off-gas como atributos adicionais. Estudos recentes demonstram que a composição dinâmica dos gases de exaustão (CO/CO$_2$) fornece informação contínua sobre o progresso da descarbonetação, permitindo predições mais precisas de $[\text{C}]_{\text{final}}$ \cite{zhou2024bof,shao2024hybrid}. A combinação desses sinais com os atributos físico-derivados estáticos propostos nesta dissertação configuraria um framework híbrido de segunda ordem.

No campo do aprendizado profundo aplicado a processos dinâmicos, a substituição do MLP estático por uma arquitetura PINN que integra as equações diferenciais de balanço de massa e energia como restrições de treinamento representa uma extensão teoricamente motivada \cite{raissi2019pinns,xia2024dmpinn}. Esta abordagem permitiria modelar a trajetória da corrida, e não apenas seu estado final, além de reduzir a dependência de grandes volumes de dados de treino.

A extensão da metodologia para o Forno Elétrico a Arco (EAF) constitui outro caminho de interesse prático e científico. O EAF compartilha com o BOF a característica de processo batelada com múltiplas variáveis de entrada e saídas de qualidade críticas, e já conta com uma base de trabalhos de modelagem híbrida que poderiam ser comparados sistematicamente com a abordagem proposta \cite{skrjanc2022eaf,logar2022eaf,manojlovic2022eaf}.

Por fim, a integração dos modelos resultantes em um gêmeo digital industrial alinha-se com as perspectivas mais recentes sobre Indústria 4.0 e sistemas ciberfísicos aplicados à siderurgia \cite{fuller2020digital,lee2015cps,hermann2016industrie4}. O teste de extrapolação (Capítulo~\ref{chp:resultados}) mostra, sob perturbação de uma variável de entrada por vez, que o ganho da engenharia de atributos físicos sobre a linha de base se mantém para $[\text{C}]_\text{final}$ e $[\text{P}]_\text{final}$; para $T_\text{final}$, nem a linha de base nem a versão físico-informada aproximam-se do limiar industrial, dentro ou fora da faixa de calibração. Como o teste cobre apenas um regime de deslocamento moderado e unidimensional, não simultâneo em múltiplas variáveis, um gêmeo digital operacional exigiria, ainda assim, monitoramento contínuo de \textit{data drift} e re-treinamento periódico, não apenas a engenharia de atributos proposta aqui.